% Schematic of the mseq_core user-facing API, shown in README.md and mseq_core/README.md.
% architecture.svg is generated from this file: edit this one, then regenerate it with
%   pdflatex architecture.tex && pdftocairo -svg architecture.pdf architecture.svg
%
% The same figure is in the MSeq NIME paper, mseq_nime2027/fig-architecture.tex: everything from
% the colour definitions down is a verbatim copy of that file, so copy it over when either changes.
% The preamble reproduces the paper's setup, nimeart's fonts (Libertine, Inconsolata) and its
% \small (8 pt on 10 pt). The picture is drawn at that size, under the paper's 17.2 cm text width,
% so it is never scaled.
%
% The style follows the KVSpan paper's figures: square boxes, straight arrows with elbows, and
% Okabe-Ito tints with darkened outlines (green for the user code, one colour per engine path,
% neutral black!4 for MIDI in and out). Frames and box outlines are thin while every connector is
% black and thicker, so the two are never confused. All text is one size; box titles are in bold,
% code identifiers are in \texttt and keep their case, every other text block starts with a
% capital.
\documentclass[tikz,border=0pt]{standalone}
\usepackage[T1]{fontenc}
\usepackage[tt=false, type1=true]{libertine}
\usepackage[varqu]{zi4}
\usepackage[libertine]{newtxmath}
\usetikzlibrary{backgrounds}
\renewcommand{\small}{\fontsize{8}{10}\selectfont}
% White background with a margin, so the figure stays readable in GitHub's dark mode
\tikzset{every picture/.append style={show background rectangle, inner frame sep=6pt,
                                      background rectangle/.style={fill=white}}}

\begin{document}
\definecolor{cG}{HTML}{009E73}  % user code: Okabe-Ito bluish green
\definecolor{cA}{HTML}{E69F00}  % direct forward: Okabe-Ito orange
\definecolor{cB}{HTML}{56B4E9}  % sequencer: Okabe-Ito sky blue
\definecolor{cC}{HTML}{9467BD}  % transport: purple
% Coordinates in cm, y pointing up, origin at the bottom left of the user code frame
\begin{tikzpicture}[
    x=1cm, y=1cm, font=\small,
    frame/.style = {draw=black!35, thin},
    % Boxes are placed by the middle of their left edge, so lines along a row stay straight
    obj/.style   = {draw=black!60, thin, fill=black!4, align=center, inner sep=1.5pt, anchor=west},
    ufrm/.style  = {frame, fill=cG!8, draw=cG!70!black},
    uobj/.style  = {obj, fill=cG!20, draw=cG!70!black},
    lane/.style  = {obj, minimum width=4.7cm},
    meth/.style  = {uobj, minimum width=3.45cm},
    ttl/.style   = {font=\small\itshape, anchor=south west, inner sep=1pt},
    tlab/.style  = {anchor=north west, inner sep=2pt},
    lbl/.style   = {inner sep=1.5pt},
    link/.style  = {line width=0.8pt},
    ar/.style    = {link, -latex},
    call/.style  = {ar, densely dashed},
    dot/.style   = {circle, fill, inner sep=0pt, minimum size=4pt},
  ]

  %% ============================== user code
  \node[ttl] at (0,4) {User code};
  \draw[frame] (0,0) rectangle (8.6,4);

  % Tracks, centred on update's row
  \draw[ufrm] (0.15,1.1) rectangle (2.85,3.15);
  \node[tlab] at (0.15,3.15) {\textbf{Tracks}};
  \node[uobj, minimum width=2.4cm, minimum height=0.5cm] at (0.3,2.5) {\textbf{\texttt{Track::play\_step}}};
  \node at (1.5,2) {$\vdots$};
  \node[uobj, minimum width=2.4cm, minimum height=0.5cm] at (0.3,1.5) {\textbf{\texttt{Track::play\_step}}};

  % Conductor
  \draw[ufrm] (4.45,0.15) rectangle (8.2,3.85);
  \node[tlab] at (4.45,3.85) {\textbf{\texttt{Conductor}}};
  \node[meth, minimum height=0.8cm] (init) at (4.6,3)
    {\textbf{\texttt{init}}\\Once, at startup};
  \node[meth, minimum height=0.8cm] (update) at (4.6,2)
    {\textbf{\texttt{update}}\\On every MIDI clock pulse};
  \node[meth, minimum height=1.1cm] (input) at (4.6,0.85)
    {\textbf{\texttt{handle\_input}}\\On every incoming message\\Channel messages only};

  \draw[call] (2.85,2) -- (update.west) node[lbl, pos=0.457, above] {Instructions};

  %% ============================== engine
  \node[ttl] at (10.35,4) {Engine};
  \draw[frame] (10.35,0) rectangle (15.35,4);

  \node[lane, minimum height=1.1cm, fill=cA!30, draw=cA!70!black] (fwd) at (10.5,3.3)
    {\textbf{Direct forward}\\\texttt{SendCC}, \texttt{SendPitchBend}, \texttt{MidiMessage}\\
     \textit{Sent right away}};
  \node[lane, minimum height=1.1cm, fill=cB!45, draw=cB!70!black] (seq) at (10.5,2)
    {\textbf{Sequencer}\\\texttt{PlayNote}, \texttt{StartNote}, \texttt{StopNote}\\
     \textit{Sent on the step}};
  \node[lane, minimum height=0.8cm, fill=cC!60, draw=cC!70!black] (tr) at (10.5,0.55)
    {\textbf{Transport}\\Clock, Start, Stop, Continue};

  % Every callback returns instructions: they merge beside the conductor, then split by kind
  \node[dot] (merge) at (8.4,2) {};
  \node[dot] (split) at (10.05,2) {};
  \draw[link] (init.east) -| (merge);
  \draw[link] (input.east) -| (merge);
  \draw[link] (update.east) -- (merge) -- (split);
  \draw[ar] (split) -- (seq.west);
  \draw[ar] (split) |- (fwd.west);

  % The callbacks drive the transport through the context
  \draw[call] (8.2,0.55) -- (tr.west);
  \node[lbl, align=center, anchor=south] at (9.475,0.6)
    {\texttt{start}, \texttt{pause},\\\texttt{resume}, \texttt{quit}};

  %% ============================== MIDI in / out
  \node[obj, minimum height=0.5cm, inner sep=2.5pt] (out) at (15.8,2) {\textbf{MIDI out}};
  \node[dot] (join) at (15.6,2) {};
  \draw[link] (fwd.east) -| (join);
  \draw[link] (tr.east) -| (join);
  \draw[link] (seq.east) -- (join);
  \draw[ar] (join) -- (out.west);

  % Channel messages go to handle_input; in slave mode, clock and transport messages drive the transport
  \node[obj, minimum height=0.5cm, inner sep=2.5pt] (in) at (0.15,-0.5) {\textbf{MIDI in}};
  \node[dot] (fork) at (input.south |- in) {};
  \draw[link] (in.east) -- (fork);
  \draw[ar] (fork) -- (input.south);
  \draw[ar] (fork) -| (tr.south);
\end{tikzpicture}
\end{document}
